\documentclass[titlepage]{article}

\usepackage[margin=1in]{geometry}
\usepackage{fancyhdr}
\usepackage{amsfonts}
\usepackage{amsmath}
\usepackage{amssymb}
\usepackage{booktabs}
\usepackage{tikz}

\title{\normalsize PLUS Session Worksheet \\[1em] \Large \textit{MATH 1554:} Linear Algebra}
\author{\textit{Review For Exam 2}}
\date{Summer 2023}

% ------- Set up header -------
\renewcommand{\sectionmark}[1]{\markright{#1}} % from https://tex.stackexchange.com/a/317393/227665
\pagestyle{fancy}
\fancyhead[L]{Summer 2023}
\fancyhead[C]{\textit{\nouppercase{\rightmark}}}
\fancyhead[R]{\textbf{\thepage}}
% -------------------------------------------------------

% ------- Set up math stuff -------
\renewcommand{\vec}[1]{\mathbf{#1}}
\DeclareMathOperator{\vspan}{Span}
\DeclareMathOperator{\row}{row}
\DeclareMathOperator{\col}{col}
\DeclareMathOperator{\colspace}{Col}
\DeclareMathOperator{\nullspace}{Nul}
\DeclareMathOperator{\rank}{rank}
% -------------------------------------------------------

% ------- Set up circle -------
\newcommand{\bubble}{\tikz[baseline=-0.625ex]\draw[radius=0.5em, fill=white] (0,0) circle;}
\newcommand{\dbubble}{\tikz[baseline=-0.625ex]\draw[radius=0.5em, fill=black] (0,0) circle;}
% -------------------------------------------------------

\begin{document}
\maketitle

\section{\centering T \& F}
\begin{table}[!h]
	\centering
	\begin{tabular}{c@{\hspace{2.5em}}cp{0.75\textwidth}}
		true & false & \\
		\specialrule{\lightrulewidth}{0.25em}{1.25em}
		\bubble & \bubble & For any \(A \in \mathbb{R}^{n \times n},\) \(A\) has \(n\) eigenvalues. \\[0.75em]
		\bubble & \bubble & For any \(A \in \mathbb{R}^{n \times n},\) \(A\) and \(A^T\) have the same eigenvalues. \\[0.75em]
		\bubble & \bubble & If \(A\vec{x} = A\vec{y}\) for any \(\vec{x}, \vec{y} \in \mathbb{R}^n\), then \(\colspace A = \{\vec{0}\}\). \\[0.75em]
		\bubble & \bubble & If \(\lambda = 0\) is not an eigenvalue for a square matrix \(A \in \mathbb{R}^{n \times n}\), 
					      then the columns of \(A\) do not form a basis for \(\mathbb{R}^n\).\\[0.75em]
		\bubble & \bubble & If \(\dim (\colspace A) = 4\) for \(A \in \mathbb{R}^{6 \times 4}\), then \(\dim (\nullspace A) = 2\). \\[0.75em]
		\bubble & \bubble & If \(A, B \in \mathbb{R}^{n \times n}\) share an eigenvector \(\vec{v} \in \mathbb{R}^n\) with a corresponding eigenvalue of 
					      \(\lambda_1, \lambda_2\), then \(\lambda_1\lambda_2\) is an eigenvalue of \(AB\). \\[0.75em]
		\bubble & \bubble & If the mapping \(\vec{x} \mapsto A\vec{x}\) scales a region in \(\mathbb{R}^2\) by a factor of 3, then \(\det A = 3\). \\[0.75em]
		\bubble & \bubble & For \(A = \begin{bmatrix}1 & 2 & 1 \\ 3 & 1 & 4\end{bmatrix}\), \(\colspace A\) is a two-dimensional subspace of 
					      \(\mathbb{R}^3\) and \(\nullspace A\) is a one-dimensional subspace of \(\mathbb{R}^2\). \\[0.75em]
		\bubble & \bubble & The number of pivots columns in \(U\) are different than \(A\) for an \(LU\)-decomposition if \(A = LU\). \\[0.75em]
		\bottomrule
	\end{tabular}
\end{table}

\section{\centering Possible/Impossible}
Indicate whether it is possible for the following statements to be true. If not, mark them as \textit{impossible}.
\begin{table}[!h]
	\centering
	\begin{tabular}{c@{\hspace{2.5em}}cp{0.75\textwidth}}
		possible & impossible & \\
		\specialrule{\lightrulewidth}{0.25em}{1.25em}
		\bubble & \bubble & \(A \in \mathbb{R}^{3 \times 3}\) with \(\nullspace A = \vspan\left\{\begin{bmatrix}1 \\ 1 \\ 0\end{bmatrix}, \begin{bmatrix}0 \\ 1 \\ 0\end{bmatrix}\right\}\) and \(\colspace A = \vspan\left\{\begin{bmatrix}2 \\ 1 \\ 0\end{bmatrix}, \begin{bmatrix}0 \\ 1 \\ 1\end{bmatrix}\right\}\).\\[2em]
		\bubble & \bubble & A matrix whose null space and column space are spanned by the same vector. \\[0.75em]
		\bubble & \bubble & \(A \in \mathbb{R}^{5 \times 4}\) with \(\dim(\colspace A) = 2 \cdot \dim(\nullspace A)\). \\[0.75em]
		\bubble & \bubble & A singular matrix with a non-zero eigenvalue. \\[0.75em]
		\bubble & \bubble & \(\lambda = 0\) is an eigenvalue of an \(n \times n\) matrix but the columns span \(\mathbb{R}^n\). \\[0.75em]
		\bottomrule
	\end{tabular}
\end{table}

\section{\centering Examples}
Provide an example of a matrix that satisfies the stated properties.
\begin{enumerate}
	\item A \(3 \times 2\) matrix whose null space is spanned by the set \(\left\{\begin{bmatrix}1 \\ 2\end{bmatrix}\right\}\) and whose column space is spanned by the set 
		\(\left\{\begin{bmatrix}1 \\ 0 \\ 2\end{bmatrix}\right\}\).
\end{enumerate}

\section{\centering Fill the Box}

\begin{enumerate}
	\item Consider the transformation that projects all vectors in \(\mathbb{R}^2\) to the line \(x_1 = x_2\). List all of the real eigenvalues of the transformation.
		\hfill\fbox{\phantom{\(\lambda_1 = 1, \lambda_2 = 0\)}}
	\item Consider the transformation that reflects all vectors in \(\mathbb{R}^2\) across the line formed by \(\vspan\left\{\begin{bmatrix}1 \\ 1\end{bmatrix}\right\}\). List all of the real 
		eigenvalues of the transformation, and an eigenvector from each respective eigenspace.

		\hfill\fbox{\phantom{\(\lambda_1 = 1, \lambda_2 = 0, \begin{bmatrix}1 \\ 1\end{bmatrix}\)}}
	\item Find all possible \(h\) such that the matrix is singular. (\textit{Hint: Use the determinant.}) \[\begin{bmatrix}1 & 2 & 3 \\ h & 4 & 1 \\ 3 & 5 & 9\end{bmatrix}\]
		\[h = \fbox{\phantom{h, h}}\]

	\item If \(\det A = 3\), and \(B\) is row-equivalent to \(A\) such that a row of \(B\) was scaled by 2, yielding the new matrix \(B'\), and two rows of \(B'\) were exchanged, yielding \(A\), 
		what is \(\det B\)? \hfill \fbox{\phantom{-3/2}}
	
	\item If \(A \in \mathbb{R}^{3 \times 4}\) with 3 pivot rows, then what is \(\dim(\nullspace A)\)? \hfill \fbox{\phantom{1}}

	\item For any stochastic matrix \(P\), what is \(\det(P - I)\)? \hfill \fbox{\phantom{0}}

	\item Let \(\mathcal{C} = \{\vec{v} \in \mathbb{R}^3 \mid A\vec{x} = \vec{v} \ \text{for some} \ \vec{x} \in \mathbb{R}^3\}\) and 
		\(A = \begin{bmatrix}1 & 0 & 2 \\ 0 & 0 & 1 \\ 1 & 1 & 1\end{bmatrix}\), what is \(\dim \mathcal{C}\)? \hfill \fbox{\phantom{23}} 

	\item Find the basis for the eigenspace corresponding to \(\lambda = 2\) for the matrix \[\begin{bmatrix}2 & 0 & 0 \\ -1 & 4 & -1 \\ -1 & 2 & 1\end{bmatrix}\] \hfill \fbox{\phantom{\(\begin{bmatrix}0 & 0 & 0 \\ 0 & 0 & 0 \\ 0 & 0 & 0\end{bmatrix}\)}}.
\end{enumerate}
\newpage
\section{Solutions}

\subsection{T \& F}
\begin{table}[!h]
	\centering
	\begin{tabular}{c@{\hspace{2.5em}}cp{0.75\textwidth}}
		true & false & \\
		\specialrule{\lightrulewidth}{0.25em}{1.25em}
		\dbubble & \bubble & For any \(A \in \mathbb{R}^{n \times n},\) \(A\) has \(n\) eigenvalues. \\[0.75em]
		\dbubble & \bubble & For any \(A \in \mathbb{R}^{n \times n},\) \(A\) and \(A^T\) have the same eigenvalues. \\[0.75em]
		\dbubble & \bubble & If \(A\vec{x} = A\vec{y}\) for any \(\vec{x}, \vec{y} \in \mathbb{R}^n\), then \(\colspace A = \{\vec{0}\}\). \\[0.75em]
		\bubble & \dbubble & If \(\lambda = 0\) is not an eigenvalue for a square matrix \(A \in \mathbb{R}^{n \times n}\), 
					      then the columns of \(A\) do not form a basis for \(\mathbb{R}^n\).\\[0.75em]
		\bubble & \dbubble & If \(\dim (\colspace A) = 4\) for \(A \in \mathbb{R}^{6 \times 4}\), then \(\dim (\nullspace A) = 2\). \\[0.75em]
		\dbubble & \bubble & If \(A, B \in \mathbb{R}^{n \times n}\) share an eigenvector \(\vec{v} \in \mathbb{R}^n\) with a corresponding eigenvalue of 
					      \(\lambda_1, \lambda_2\), then \(\lambda_1\lambda_2\) is an eigenvalue of \(AB\). \\[0.75em]
		\bubble & \dbubble & If the mapping \(\vec{x} \mapsto A\vec{x}\) scales a region in \(\mathbb{R}^2\) by a factor of 3, then \(\det A = 3\). \\[0.75em]
		\bubble & \dbubble & For \(A = \begin{bmatrix}1 & 2 & 1 \\ 3 & 1 & 4\end{bmatrix}\), \(\colspace A\) is a two-dimensional subspace of 
					      \(\mathbb{R}^3\) and \(\nullspace A\) is a one-dimensional subspace of \(\mathbb{R}^2\). \\[0.75em]
		\bubble & \dbubble & The number of pivots columns in \(U\) are different than \(A\) for an \(LU\)-decomposition if \(A = LU\). \\[0.75em]
		\bottomrule
	\end{tabular}
\end{table}

\subsection{Possible/Impossible}
\begin{table}[!h]
	\centering
	\begin{tabular}{c@{\hspace{2.5em}}cp{0.75\textwidth}}
		possible & impossible & \\
		\specialrule{\lightrulewidth}{0.25em}{1.25em}
		\bubble & \dbubble & \(A \in \mathbb{R}^{3 \times 3}\) with \(\nullspace A = \vspan\left\{\begin{bmatrix}1 \\ 1 \\ 0\end{bmatrix}, \begin{bmatrix}0 \\ 1 \\ 0\end{bmatrix}\right\}\) and \(\colspace A = \vspan\left\{\begin{bmatrix}2 \\ 1 \\ 0\end{bmatrix}, \begin{bmatrix}0 \\ 1 \\ 1\end{bmatrix}\right\}\).\\[2em]
		\dbubble & \bubble & A matrix whose null space and column space are spanned by the same vector. \\[0.75em]
		\bubble & \dbubble & \(A \in \mathbb{R}^{5 \times 4}\) with \(\dim(\colspace A) = 2 \cdot \dim(\nullspace A)\). \\[0.75em]
		\dbubble & \bubble & A singular matrix with a non-zero eigenvalue. \\[0.75em]
		\bubble & \dbubble & \(\lambda = 0\) is an eigenvalue of an \(n \times n\) matrix but the columns span \(\mathbb{R}^n\). \\[0.75em]
		\bottomrule
	\end{tabular}
\end{table}

\subsection{Examples}
\begin{enumerate}
	\item There are multiple answers. The simplest is \[\begin{bmatrix}2 & -1 \\ 0 & 0 \\ 4 & -2\end{bmatrix}.\] In general, if the first column is \(c\begin{bmatrix}1 \\ 0 \\ 2\end{bmatrix}\), 
		then the second column must be \(-\frac{c}{2}\begin{bmatrix}1 \\ 0 \\ 2\end{bmatrix}\). In this case, I chose \(c = 2\).
\end{enumerate}

\subsection{Fill in the Box}
\begin{enumerate}
	\item \(\lambda_1 = 1, \lambda_2 = 0\).
	\item \(\lambda_1 = 1, \lambda_2 = -1\) with a corresponding \(\vec{v}_1 = \begin{bmatrix}1 \\ 1\end{bmatrix}\) and \(\vec{v}_2 = \begin{bmatrix}-1 \\ 1\end{bmatrix}\).
	\item \(h = 1/3\). If we chose to perform co-factor expansion on the first column, then we get 
		\[\begin{vmatrix}1 & 2 & 3 \\ h & 4 & 1 \\ 3 & 5 & 9\end{vmatrix} 
		  = 1\begin{vmatrix}4 & 1 \\ 5 & 9\end{vmatrix} - h\begin{vmatrix}2 & 3 \\ 5 & 9\end{vmatrix} + 3 \begin{vmatrix}2 & 3 \\ 4 & 1\end{vmatrix} = 1(31) - h(3) + 3(-10) = -3h + 1\]
		We need the matrix to be \emph{singular}, which implies that is has to have a determinant of 0. We've got an expression for the determinant, and now we need a value of \(h\) that
		makes the determinant equal to 0. \[-3h + 1 = 0 \implies h = 1/3\].
	\item \(\det B = -\frac{3}{2}\). Informally,		
		\[\det B \overset{2 \cdot \row B}{\longrightarrow} 2\det B \overset{\text{Row swap} \ B'}{\longrightarrow} -2 \det B = \det A = 3.\]
		Meaning, \[-2\det B = 3 \implies \det B = -\frac{3}{2}\]
	\item \textbf{1}.
	\item \textbf{0}.
	\item \textbf{3}.
	\item We need to find a basis for \(\nullspace (A - 2I)\). 
		\[A - 2I = \begin{bmatrix}0 & 0 & 0 \\ -1 & 2 & -1 \\ -1 & 2 & -1\end{bmatrix} \sim \begin{bmatrix}1 & -2 & 1 \\ 0 & 0 & 0 \\ 0 & 0 & 0\end{bmatrix}\]
		From here, we deduce that \(\vec{x} = \begin{bmatrix}2x_2 - x_3 \\ x_2 \\ x_3\end{bmatrix} = x_2\begin{bmatrix}2 \\ 1\\ 0\end{bmatrix} + x_3\begin{bmatrix}-1 \\ 0\\ 1\end{bmatrix}\). So a basis would be \[\left\{\begin{bmatrix}2 \\ 1\\ 0\end{bmatrix}, \begin{bmatrix}-1 \\ 0\\ 1\end{bmatrix}\right\}.\]
\end{enumerate}

\end{document}